\section{Introduction}
\label{sec:intro}

As machine learning progresses rapidly, its societal impact has come under
scrutiny. An important concern is potential discrimination based on protected
attributes such as gender, race, or religion. Since learned predictors and risk
scores increasingly support or even replace human judgment, there is an
opportunity to formalize what harmful discrimination means and to design
algorithms that avoid it. However, researchers have found it difficult to agree
on a single measure of discrimination. As of now, there are several competing
approaches, representing different opinions and striking different trade-offs.
Most of the proposed fairness criteria are observational: They depend only on
the joint distribution of predictor~$\hY,$ protected attribute~$A$,
features~$X$, and outcome $Y.$ For example, the natural requirement that~$\hY$
and~$A$ must be statistically independent is referred to as \emph{demographic
parity}. Some approaches transform the features~$X$ to obfuscate the information
they contain about~$A$~\cite{Zemel2013}. The recently proposed \emph{equalized
odds} constraint~\cite{Hardt2016} demands that the predictor $\hY$ and the
attribute $A$ be independent conditional on the actual outcome $Y.$ All three
are examples of observational approaches.

A growing line of work points at the insufficiency of existing definitions.
Hardt, Price and Srebro~\cite{Hardt2016} construct two scenarios with
\emph{intuitively} different social interpretations that admit identical joint
distributions over $(\hY, A, Y, X)$. Thus, no observational criterion can
distinguish them. While there are non-observational criteria, notably the early
work on individual fairness~\cite{Dwork2012}, these have not yet gained
traction.  So, it might appear that the community has reached an impasse.

% ----------------------------- SUBSECTION  -----------------------------------
\subsection{Our contributions}

We assay the problem of discrimination in machine learning in the language of
causal reasoning. This viewpoint supports several contributions:

\begin{itemize}
  \item Revisiting the two scenarios proposed in~\cite{Hardt2016}, we articulate
    a natural causal criterion that formally distinguishes them. In particular,
    we show that observational criteria are unable to determine if a protected
    attribute has \emph{direct causal influence} on the predictor that is not
    mitigated by \emph{resolving} variables.
  \item We point out subtleties in fair decision making that arise naturally
    from a causal perspective, but have gone widely overlooked in the past.
    Specifically, we formally argue for the need to distinguish between the
    underlying concept behind a protected attribute, such as race or gender, and
    its \emph{proxies} available to the algorithm, such as visual features or
    name.
  \item We introduce and discuss two natural causal criteria centered around the
    notion of \emph{interventions} (relative to a causal graph) to formally
    describe specific forms of discrimination.
  \item Finally, we initiate the study of algorithms that avoid these forms of
    discrimination. Under certain linearity assumptions about the underlying
    causal model generating the data, an algorithm to remove a specific kind of
    discrimination leads to a simple and natural heuristic.
\end{itemize}

At a higher level, our work proposes a shift from trying to find a single
statistical fairness criterion to arguing about properties of the data and which
assumptions about the generating process are justified. Causality provides a
flexible framework for organizing such assumptions.

% ----------------------------- SUBSECTION  -----------------------------------
\subsection{Related work}

Demographic parity and its variants have been discussed in numerous papers,
e.g., \cite{Feldman2015, Zemel2013, Zafar2017, Edwards2015}. While demographic
parity is easy to work with, the authors of \cite{Dwork2012} already highlighted
its insufficiency as a fairness constraint. In an attempt to remedy the
shortcomings of demographic parity~\cite{Hardt2016} proposed two notions,
\emph{equal opportunity} and \emph{equal odds}, that were also considered
in~\cite{Zafar2017a}.  A review of various fairness criteria can be found
in~\cite{Berk2017}, where they are discussed in the context of criminal justice.
In~\cite{Kleinberg2016, Chouldechova2016} it has been shown that imperfect
predictors cannot simultaneously satisfy equal odds and \emph{calibration}
unless the groups have identical base rates, \ie{} rates of positive outcomes.

A starting point for our investigation is the unidentifiability result
of~\cite{Hardt2016}. It shows that observedvational criteria are too weak to
distinguish two intuitively very different scenarios. However, the work does not
provide a formal mechanism to articulate why and how these scenarios should be
considered different. Inspired by Pearl's causal interpretation of Simpson's
paradox~\cite[Section 6]{Pearl2009}, we propose causality as a way of coping
with this unidentifiability result.

An interesting non-observational fairness definition is the notion of
\emph{individual fairness}~\cite{Dwork2012} that assumes the existence of a
similarity measure on individuals, and requires that any two similar individuals
should receive a similar distribution over outcomes. More recent work lends
additional support to such a definition~\cite{Friedler2016}. From the
perspective of causality, the idea of a similarity measure is akin to the method
of \emph{matching} in counterfactual reasoning~\cite{Rosenbaum1983,
Qureshi2016}. That is, evaluating approximate counterfactuals by comparing
individuals with similar values of covariates excluding the protected attribute.

Recently,~\cite{Kusner2017} put forward one possible causal definition, namely
the notion of \emph{counterfactual fairness}. It requires modeling
counterfactuals on a per individual level, which is a delicate task. Even
determining the effect of \emph{race} at the group level is difficult; see the
discussion in~\cite{VanderWeele2014}. The goal of our paper is to assay a more
general causal framework for reasoning about discrimination in machine learning
without committing to a single fairness criterion, and without committing to
evaluating individual causal effects. In particular, we draw an explicit
distinction between the protected attribute (for which interventions are often
impossible in practice) and its proxies (which sometimes can be intervened
upon).

Moreover, causality has already been employed for the discovery of
discrimination in existing data sets by~\cite{Bonchi2015, Qureshi2016}. Causal
graphical conditions to identify \emph{meaningful partitions} have been proposed
for the discovery and prevention of certain types of discrimination by
preprocessing the data~\cite{Zhang2017b}. These conditions rely on the
evaluation of \emph{path specific effects}, which can be traced back all the way
to~\cite[Section 4.5.3]{Pearl2009}. The authors of~\cite{Nabi2017} recently
picked up this notion and generalized Pearl's approach by a constraint based
prevention of discriminatory path specific effects arising from counterfactual
reasoning. Our research was done independently of these works.

% ----------------------------- SUBSECTION  -----------------------------------
\subsection{Causal graphs and notation}

Causal graphs are a convenient way of organizing assumptions about the data
generating process. We will generally consider causal graphs involving a
protected attribute~$A,$ a set of proxy variables~$P,$ features~$X,$ a
predictor~$R$ and sometimes an observed outcome~$Y.$ For background on causal
graphs see~\cite{Pearl2009}. In the present paper a \emph{causal graph} is a
directed, acyclic graph whose nodes represent random variables. A \emph{directed
path} is a sequence of distinct nodes~$V_1, \dots, V_k$, for~$k \ge 2$, such
that~$V_i \to V_{i+1}$ for all~$i \in \{1, \dots,k-1\}$. We say a directed path
is \emph{blocked by a set of nodes~$Z$}, where~$V_1, V_k \notin Z$, if~$V_i \in
Z$ for some~$i \in \{2, \dots, k-1 \}$.\footnote{As it is not needed in our
work, we do not discuss the graph-theoretic notion of d-separation.}

A \emph{structural equation model} is a set of equations~$V_i = f_i(pa(V_i),
N_i)$, for~$i \in \{1, \dots, n\}$, where~$pa(V_i)$ are the parents of~$V_i$,
\ie{} its \emph{direct causes}, and the~$N_i$ are independent noise variables. We interpret
these equations as assignments. Because we assume acyclicity, starting from the
roots of the graph, we can recursively compute the other variables, given the
noise variables. This leads us to view the structural equation model and its
corresponding graph as a \emph{data generating model}. The predictor~$R$ maps
inputs, e.g., the features~$X$, to a predicted output. Hence we model it as a
childless node, whose parents are its input variables. Finally, note that given
the noise variables, a structural equation model entails a unique joint
distribution; however, the same joint distribution can usually be entailed by
multiple structural equation models corresponding to distinct causal structures.
